\chapter{諸定理}
    \section{偶関数と奇関数への分解}
        \begin{shadebox}
            連続な1変数関数$f(x)$は連続な偶関数,奇関数に一通りに分解できる
        \end{shadebox}
        \def\fe{f_\mathrm{e}}
        \def\fo{f_\mathrm{o}}
        \def\fea{f_\mathrm{ea}}
        \def\foa{f_\mathrm{oa}}
        \def\feb{f_\mathrm{eb}}
        \def\fob{f_\mathrm{ob}}
        \begin{proof}
            \quad\par
            (分解可能性)\\
            \begin{align*}
                \fe(x) &\coloneqq \frac{1}{2}\left(f(x)+f(-x)\right)\\
                \fo(x) &\coloneqq \frac{1}{2}\left(f(x)-f(-x)\right)
            \end{align*}
            とすれば確かに$f(x)=\fe(x)+\fo(x)$であり$\fe(x),\fo(x)$はそれぞれ連続な偶,奇関数になっている。\\
            \newline
            (一意性)
            \par
            次に、分解が一通りであることを示す。$f(x)$の分解が次のように2通りあったとする。
            \begin{align}
                \begin{cases}
                f(x) = \fea(x) + \foa(x) & \\
                f(x) = \feb(x) + \fob(x) &
                \end{cases}
            \end{align}
            任意の$l>0$に対して
            \begin{align*}
                0 &= \integrate{-l}{l}{\left(f(x)-f(x)\right)}{}{x} = \integrate{-l}{l}{\left(\fea(x) + \foa(x) - \feb(x) + \fob(x)\right)}{}{x} \\
                &= 2\integrate{0}{l}{\left(\fea(x)-\feb(x)\right)}{}{x}
            \end{align*}
            であるから両辺を$l$で微分すると
            \[0 = \fea(l) - \feb(l)\]
            これで$x>0$において$\fea,\feb$が等しいことが示されたが、そもそも両者は偶関数であったから$x<0$においても等しくならざるを得ない。
            残るは$x=0$であるが、$\foa,\fob$は連続な奇関数なので$\foa(0)=\fob(0)=0$。
            これと式(1)より$ \fea(0) = \feb(0) = f(0)$となり、結局$\fea$と$\feb$は常に等しい。
            これと式(1)から奇関数成分$\foa,\fob$も等しいことがわかる。
        \end{proof}
    \section{\texorpdfstring{$\sup_{\bm{x}\in X}(f(\bm{x})+g(\bm{x})) \leq \sup_{\bm{x}\in X}f(\bm{x}) + \sup_{\bm{x}\in X}g(\bm{x}),\;\inf_{\bm{x}\in X}f(\bm{x}) + \inf_{\bm{x}\in X}g(\bm{x}) \leq \inf_{\bm{x}\in X}(f(\bm{x})+g(\bm{x}))$}{sup\_\{x∈X\} f(x)+g(x) ≦ sum\_\{x∈X\}f(x) + sup\_\{x∈X\}g(x), inf\_\{x∈X\}f(x) + inf\_\{x∈X\}g(x) ≦ inf\_\{x∈X\} f(x)+g(x)}}
        \begin{shadebox}
            集合$X$上で定義された有界な関数$f,g$について次が成り立つ。
            \[ \sup_{\bm{x}\in X}(f(\bm{x})+g(\bm{x})) \leq \sup_{\bm{x}\in X}f(\bm{x}) + \sup_{\bm{x}\in X}g(\bm{x}),\;\inf_{\bm{x}\in X}f(\bm{x}) + \inf_{\bm{x}\in X}g(\bm{x}) \leq \inf_{\bm{x}\in X}(f(\bm{x})+g(\bm{x})) \]
        \end{shadebox}
        \begin{proof}
            \quad\par
            $\sup_{\bm{x}\in X}(f(\bm{x})+g(\bm{x}))$を与える$\bm{x}\in X$の一つを$\bm{x}_1$とすると$\sup_{\bm{x}\in X}(f(\bm{x})+g(\bm{x})) = f(\bm{x}_1)+g(\bm{x}_1) \leq \sup_{\bm{x}\in X}f(\bm{x}) + \sup_{\bm{x}\in X}g(\bm{x})$。
            また、$\inf_{\bm{x}\in X}(f(\bm{x})+g(\bm{x}))$を与える$\bm{x}\in X$の一つを$\bm{x}_2$とすると$\inf_{\bm{x}\in X}(f(\bm{x})+g(\bm{x})) = f(\bm{x}_2)+g(\bm{x}_2) \geq \inf_{\bm{x}\in X}f(\bm{x}) + \inf_{\bm{x}\in X}g(\bm{x})$
        \end{proof}
    \section{\texorpdfstring{$\sup_{\bm{x}\in X}(-f(\bm{x})) = -\inf_{\bm{x}\in X}f(\bm{x}),\; \inf_{\bm{x}\in X}(-f(\bm{x})) = -\sup_{\bm{x}\in X}f(\bm{x})$}{sup\_\{x∈X\} -f(x) = -inf\_\{x∈X\}f(x), inf\_\{x∈X\} -f(x) = -sup\_\{x∈X\}f(x)}}
        \begin{shadebox}
            集合$X$上で定義された有界な関数$f,g$について次が成り立つ。
            \[ \sup_{\bm{x}\in X}(-f(\bm{x})) = -\inf_{\bm{x}\in X}f(\bm{x}),\; \inf_{\bm{x}\in X}(-f(\bm{x})) = -\sup_{\bm{x}\in X}f(\bm{x}) \]
        \end{shadebox}
        \begin{proof}
            \quad\par
            $^\forall\bm{y}\in X,\;\inf_{\bm{x}\in X}f(\bm{x}) \leq f(\bm{y})$なので$^\forall\bm{y}\in X,\;-\inf_{\bm{x}\in X}f(\bm{x}) \geq -f(\bm{y})$。
            任意の$\varepsilon>0$に対して$\inf_{\bm{x}\in X}f(\bm{x}) + \varepsilon$は$f$の下界ではないから、適当な$\bm{x}_1\in X$が存在して$f(\bm{x}_1) < \inf_{\bm{x}\in X}f(\bm{x}) + \varepsilon \quad \therefore\; -f(\bm{x}_1) > -\inf_{\bm{x}\in X}f(\bm{x}) - \varepsilon$
            。すなわち$-\inf_{\bm{x}\in X}f(\bm{x}) - \varepsilon$は$-f$の上界ではない。
            よって上限の定義より$\sup_{\bm{x}\in X}(-f(\bm{x})) = -\inf_{\bm{x}\in X}f(\bm{x})$。2つ目の主張の証明も同様。
        \end{proof}
    \section{H\"older(ヘルダー)の不等式}
        \begin{shadebox}
            $p,q>1,\;\frac{1}{p}+\frac{1}{q}=1$のとき
            \[\sum_{i=1}^n|x_iy_i|\leq\norm{\bm{x}}_p\norm{\bm{y}}_q\]
        \end{shadebox}
        \begin{proof}
            \quad\par
            $\bm{x},\bm{y}$の少なくとも1つが零ベクトルであるときは明らかに成り立つ。
            以下では両者とも零ベクトルでないとする。
            Youngの不等式より$a,b\geq 0$に対して$ab \leq a^p/p + b^q/q$が成り立つ。
            この式で$a=|x_i|/\norm{\bm{x}}_p,b=|y_i|/\norm{\bm{y}}_q$とすると
            \[\frac{|x_iy_i|}{\norm{\bm{x}}_p\norm{\bm{y}}_q}\leq\frac{|x_i|^p/\norm{\bm{x}}_p^p}{p}+\frac{|y_i|^q/\norm{\bm{y}}_q^q}{q}\]
            これを$i=1,\dots,n$まで足し合わせると
            \[\frac{1}{\norm{\bm{x}}_p\norm{\bm{y}}_q}\sum_{i=1}^n{|x_iy_i|} \leq \frac{\norm{\bm{x}}_p^p/\norm{\bm{x}}_p^p}{p}+\frac{\norm{\bm{y}}_q^q/\norm{\bm{y}}_q^q}{q} = 1/p+1/q = 1\]
            \[\therefore\;\sum_{i=1}^n|x_iy_i|\leq\norm{\bm{x}}_p\norm{\bm{y}}_q\]
        \end{proof}
    \section{Minkowskiの不等式(三角不等式の一般化)}
        \begin{shadebox}
            $1<p<\infty$のとき
            \[\norm{\bm{x}+\bm{y}}_p \leq \norm{\bm{x}}_p+\norm{\bm{y}}_p\]
        \end{shadebox}
        \begin{proof}
            \begin{equation}
                \norm{\bm{x}+\bm{y}}_p^p = \sum_{i=1}^n{|x_i+y_i|^p} = \sum_{i=1}^n{|x_i+y_i||x_i+y_i|^{p-1}} \leq \sum_{i=1}^n{|x_i||x_i+y_i|^{p-1}}+\sum_{i=1}^n{|y_i||x_i+y_i|^{p-1}}
            \end{equation}
            右辺第1項はH\"olderの不等式より上から
            \[\norm{\bm{x}}_p\norm{\bm{v}}_{\frac{p}{p-1}}\quad\mathrm{where}\quad\bm{v}[i]\coloneqq|x_i+y_i|^{p-1}\]
            で抑えられる。さらに
            \[\norm{\bm{x}}_p\norm{\bm{v}}_{\frac{p}{p-1}} = \norm{\bm{x}}_p\left[\sum_{i=1}^n{|x_i+y_i|^{(p-1)\frac{p}{p-1}}}\right]^{\frac{p-1}{p}} = \norm{\bm{x}}_p\norm{\bm{x}+\bm{y}}_p^{p-1}\]
            同様に式(1)の右辺第2項も上から$\norm{\bm{y}}_p\norm{\bm{x}+\bm{y}}_p^{p-1}$で抑えられるので
            \[\norm{\bm{x}+\bm{y}}_p^p \leq \norm{\bm{x}}_p\norm{\bm{x}+\bm{y}}_p^{p-1} + \norm{\bm{y}}_p\norm{\bm{x}+\bm{y}}_p^{p-1}\]
            両辺を$\norm{\bm{x}+\bm{y}}_p^{p-1}$で割ることでMinkowskiの不等式を得る。
        \end{proof}